\documentclass[11pt]{article}
\usepackage[a4paper,margin=25mm]{geometry}
\usepackage{fontspec,polyglossia}
\defaultfontfeatures{Extension=.otf}
\setmainfont{NimbusRoman}[UprightFont=*-Regular,BoldFont=*-Bold,ItalicFont=*-Italic,BoldItalicFont=*-BoldItalic]
\setsansfont{NimbusSans}[UprightFont=*-Regular,BoldFont=*-Bold,ItalicFont=*-Italic,BoldItalicFont=*-BoldItalic]
\setmonofont{NimbusMonoPS}[UprightFont=*-Regular,BoldFont=*-Bold,ItalicFont=*-Italic,BoldItalicFont=*-BoldItalic,Scale=0.85]
\newfontfamily\cyrillicfont{NimbusRoman}[UprightFont=*-Regular,BoldFont=*-Bold,ItalicFont=*-Italic,BoldItalicFont=*-BoldItalic,Script=Cyrillic]
\newfontfamily\cyrillicfontsf{NimbusSans}[UprightFont=*-Regular,BoldFont=*-Bold,ItalicFont=*-Italic,BoldItalicFont=*-BoldItalic,Script=Cyrillic]
\newfontfamily\cyrillicfonttt{NimbusMonoPS}[UprightFont=*-Regular,BoldFont=*-Bold,ItalicFont=*-Italic,BoldItalicFont=*-BoldItalic,Script=Cyrillic,Scale=0.85]
\def\ruCode{ru}
\ifx\wrehLang\ruCode
  \setdefaultlanguage{russian}\setotherlanguage{english}
  \newcommand{\Lang}[2]{#2}
\else
  \setdefaultlanguage{english}\setotherlanguage{russian}
  \newcommand{\Lang}[2]{#1}
\fi
\usepackage{amsmath,amssymb,amsthm,mathtools}
\usepackage{microtype,booktabs,tabularx,longtable,array,graphicx}
\newcolumntype{Y}{>{\raggedright\arraybackslash}X}
\usepackage{enumitem,xcolor}
\usepackage[hidelinks,unicode]{hyperref}
\definecolor{wrehblue}{RGB}{22,71,105}
\newtheorem{definition}{\Lang{Definition}{Определение}}[section]
\newtheorem{proposition}[definition]{\Lang{Proposition}{Утверждение}}
\newtheorem{theorem}[definition]{\Lang{Theorem}{Теорема}}
\newtheorem{corollary}[definition]{\Lang{Corollary}{Следствие}}
\theoremstyle{definition}
\newtheorem{example}[definition]{\Lang{Example}{Пример}}
\newtheorem{remark}[definition]{\Lang{Remark}{Замечание}}
\newcommand{\Wcal}{\mathcal W}
\newcommand{\Ccal}{\mathfrak C}
\newcommand{\Hminus}{\mathcal H_-}
\newcommand{\Pcal}{\mathfrak H_-}
\newcommand{\Ocal}{\mathfrak O_-}
\newcommand{\E}{\mathbb E}
\newcommand{\Var}{\operatorname{Var}}
\newcommand{\Cov}{\operatorname{Cov}}
\newcommand{\supp}{\operatorname{supp}}
\newcommand{\figfile}[1]{figures/#1_\wrehLang.pdf}
\setlength{\emergencystretch}{2em}
\setlist[enumerate]{itemsep=3pt,topsep=4pt}
\renewcommand{\arraystretch}{1.18}
\hypersetup{
 pdftitle={\Lang{WR-IV: Response-Conditioned Global Completion and the Status of the Past}{WR-IV: Глобальное достраивание при заданных откликах и статус прошлого}},
 pdfauthor={A. A. Malachevsky},
 pdfsubject={WREH; bilingual technical preprint v0.6; unpeer-reviewed},
 pdfkeywords={historical projection, global realizability, Brownian bridge, conditional distribution, smoothing, epistemic horizons}
}
\title{\color{wrehblue}\textbf{WR-IV}\\[5pt]
\Lang{Response-Conditioned Global Completion\\and the Status of the Past}{Глобальное достраивание при заданных откликах\\и статус прошлого}\\[9pt]
\large\Lang{Technical preprint v0.6}{Технический препринт v0.6}}
\author{A. A. Malachevsky\\
\small\href{https://orcid.org/0009-0008-6009-3196}{ORCID: 0009-0008-6009-3196}\\
\small World Realizability \& Epistemic Horizons (WREH)}
\date{\Lang{8 October 2026}{8 октября 2026 года}}
